\SourceLeaf{331}

oder

\begin{equation}
\sin 1 \, . \, \sin(2+3)\sin(2+5) = \sin 2 \, . \, \sin(1+4)\sin(1+6).
\tag{3}
\end{equation}

Auf gleiche Weise oder auch bloss durch Permutation

\begin{equation}
\sin 1 \, . \, \sin(4+3)\sin(4+5) = \sin 4 \, . \, \sin(1+2)\sin(1+6)
\tag{4}
\end{equation}

\begin{equation}
\sin 1 \, . \, \sin(6+3)\sin(6+5) = \sin 6 \, . \, \sin(1+2)\sin(1+4)
\tag{5}
\end{equation}

\begin{equation}
\sin 3 \, . \, \sin(2+1)\sin(2+5) = \sin 2 \, . \, \sin(3+4)\sin(3+6)
\tag{6}
\end{equation}

\begin{equation}
\sin 3 \, . \, \sin(4+1)\sin(4+5) = \sin 4 \, . \, \sin(3+2)\sin(3+6)
\tag{7}
\end{equation}

\begin{equation}
\sin 3 \, . \, \sin(6+1)\sin(6+5) = \sin 6 \, . \, \sin(3+2)\sin(3+4)
\tag{8}
\end{equation}

\begin{equation}
\sin 5 \, . \, \sin(2+1)\sin(2+3) = \sin 2 \, . \, \sin(5+4)\sin(5+6)
\tag{9}
\end{equation}

\begin{equation}
\sin 5 \, . \, \sin(4+1)\sin(4+3) = \sin 4 \, . \, \sin(5+2)\sin(5+6)
\tag{10}
\end{equation}

\begin{equation}
\sin 5 \, . \, \sin(6+1)\sin(6+3) = \sin 6 \, . \, \sin(5+2)\sin(5+4).
\tag{11}
\end{equation}

\Needspace{5\baselineskip}Diese neun Gleichungen können auch so geschrieben werden:

\begin{equation}
\sin 1 \, . \, \sin\mathfrak{C} \, \sin(\mathfrak{B}-B) = \sin 2 \, . \, \sin(\mathfrak{C}-C)\sin\mathfrak{B}
\tag{3}
\end{equation}

\begin{equation}
\sin 1 \, . \, \sin B \, \sin\mathfrak{A} = \sin 4 \, . \, \sin A \, \sin\mathfrak{B}
\tag{4}
\end{equation}

\begin{equation}
\sin 1 \, . \, \sin(\mathfrak{A}-A)\sin C = \sin 6 \, . \, \sin A \, \sin(\mathfrak{C}-C)
\tag{5}
\end{equation}

\begin{equation}
\sin 3 \, . \, \sin A \, \sin(\mathfrak{B}-B) = \sin 2 \, . \, \sin B \, \sin(\mathfrak{A}-A)
\tag{6}
\end{equation}

\begin{equation}
\sin 3 \, . \, \sin(\mathfrak{C}-C)\sin\mathfrak{A} = \sin 4 \, . \, \sin\mathfrak{C} \, \sin(\mathfrak{A}-A)
\tag{7}
\end{equation}

\begin{equation}
\sin 3 \, . \, \sin\mathfrak{B} \, \sin C = \sin 6 \, . \, \sin\mathfrak{C} \, \sin B
\tag{8}
\end{equation}

\begin{equation}
\sin 5 \, . \, \sin A \, \sin\mathfrak{C} = \sin 2 \, . \, \sin\mathfrak{A} \, \sin C
\tag{9}
\end{equation}

\begin{equation}
\sin 5 \, . \, \sin(\mathfrak{C}-C)\sin B = \sin 4 \, . \, \sin(\mathfrak{B}-B)\sin C
\tag{10}
\end{equation}

\begin{equation}
\sin 5 \, . \, \sin\mathfrak{B} \, \sin(\mathfrak{A}-A) = \sin 6 \, . \, \sin(\mathfrak{B}-B)\sin\mathfrak{A},
\tag{11}
\end{equation}

indem

\[
A=1+2, \qquad B=3+4, \qquad C=5+6,
\]
\[
180^0-\mathfrak{A}=4+5, \qquad 180^0-\mathfrak{B}=6+1, \qquad 180^0-\mathfrak{C}=2+3,
\]
\[
\mathfrak{A}-A=3+6, \qquad \mathfrak{B}-B=2+5, \qquad \mathfrak{C}-C=1+4
\]
ist.

\Needspace{9\baselineskip}Die Ableitung kann auch aus Betrachtung der Figur geschehen, nemlich
\begin{center}(1) und (2) nach Gefallen aus I oder II\end{center}
